\documentclass[12pt]{article}
\usepackage{fullpage}
\usepackage{amsmath,amssymb,amsthm}

\newcommand{\cL}{\mathcal L}
\newcommand{\Dr}{\operatorname{Dr}}
\newcommand{\Trop}{\operatorname{Trop}}
\newcommand{\rk}{\operatorname{rk}}
\newcommand{\cl}{\operatorname{cl}}
\newcommand{\op}{\operatorname{op}}

\newtheorem{theorem}{Theorem}
\newtheorem{lemma}{Lemma}
\newtheorem{proposition}{Proposition}

\begin{document}

\begin{center}
{\bf A canonical opposition map on the relative Dressian}
\end{center}

\section*{Conventions and standard inputs}

Let $E$ be the ground set, let $r=\rk M$, and let $\mu_M$ be the
trivial rank $r$ valuated matroid supported on the bases of $M$.  I
write
\[
        p\preceq q
\]
to mean that the valuated matroid $p$ is a valuated quotient of $q$.
With the convention of Brandt--Eur--Zhang, a rank $k$ point of
$\Dr(M)$ is a rank $k$ valuated matroid $\nu$ with
$\nu\preceq\mu_M$, and
\[
        \Trop(\theta)\subseteq \Trop(\nu)
        \quad\Longleftrightarrow\quad
        \theta\preceq\nu .
\]
Valuated matroids are projective: adding one scalar to all finite
Plucker coordinates does not change the tropical linear space.

For ranks $a\le b$, the quotient relation is characterized by the
tropical incidence-Plucker relations
\[
 \min_{i\in T\setminus S}\{p(Si)+q(T-i)\}
 \quad\hbox{is attained at least twice}                    \tag{1}
\]
for $|S|=a-1$, $|T|=b+1$, $S\subseteq T$, together with the valuated
matroid axioms.  For $b=a+1$ this is the three-term relation
\[
 \min\{p(Sx)+q(Syz),\,p(Sy)+q(Sxz),\,p(Sz)+q(Sxy)\}.
                                                               \tag{2}
\]
I use the following standard consequences of the Dress--Wenzel--Murota
cryptomorphisms.  First, a complete sequence
$g_0,\ldots,g_r$ with nonempty layers, $g_0$ the trivial rank $0$
valuation, and $g_r=\mu_M$ is a valuated complete flag if and only if
all adjacent relations (2) hold.  This is the local
$M^\natural$-concavity criterion; see Murota, Theorem 3.5 and
Corollary 3.4, and, for the flag interpretation, Brandt--Eur--Zhang,
Proposition 4.2.3 and Theorem 4.3.1, and Jarra--Lorscheid,
Theorem 2.17 and Proposition 2.21.  In the finite uniform adjacent-rank
case, Joswig--Loho--Luber--Olarte, Theorem 3.10, gives a direct
incidence-implies-Plucker statement.  Secondly, a single quotient
$\nu\preceq\mu_M$ of rank $k$ can be inserted into such a complete
flag by valuated truncation and relative Higgs elongation:
\[
 h_j(B)=
 \begin{cases}
 \min\{\nu(C): B\subseteq C,\ |C|=k\}, & j\le k,\ |B|=j,\\
 \min\{\nu(C): C\subseteq B,\ |C|=k\}, & j\ge k,\ B\in\mathcal I_j(M),\\
 \infty, & j\ge k,\ B\notin\mathcal I_j(M).
 \end{cases}                                                \tag{3}
\]
After harmless scalar normalizations, these layers satisfy
$h_0\preceq\cdots\preceq h_r=\mu_M$ and $h_k=\nu$.  This is Murota's
valuated truncation/elongation (Theorem 2.1 in the 1997 paper),
equivalently the $M^\natural$-convex infimal convolution construction
in discrete convex analysis.

\section*{1. Flat-lattice consequences}

\begin{lemma}\label{lem:modular}
If $L=\cL(M)$ admits an order-reversing involution
$x\mapsto x^\perp$, then $L$ is modular and
$\rk(x^\perp)=r-\rk(x)$.
\end{lemma}

\begin{proof}
The flat lattice is a finite ranked semimodular lattice.  An
anti-automorphism reverses maximal chains, hence sends rank $i$
elements to rank $r-i$ elements.  Applying the anti-automorphism to
the semimodular inequality gives the reverse inequality, so equality
holds for every pair.  This is modularity.
\end{proof}

Thus $(x\vee y)^\perp=x^\perp\wedge y^\perp$ and
$(x\wedge y)^\perp=x^\perp\vee y^\perp$.

\begin{lemma}[flat constancy]\label{lem:flatconst}
Let $\nu\preceq \mu_M$ have rank $k$.  If $B,B'$ are independent
$k$-subsets of $M$ with $\cl(B)=\cl(B')$, then $\nu(B)=\nu(B')$.
Moreover $\nu(A)=\infty$ whenever $A$ is dependent in $M$.
\end{lemma}

\begin{proof}
The underlying matroid of $\nu$ is a quotient of $M$, so its bases are
independent in $M$; this gives the support assertion.

For $0<k\le r$, the basis graph of the restriction to a fixed rank
$k$ flat is connected.  Hence it suffices to compare
$B=A\cup\{b\}$ and $B'=A\cup\{b'\}$ with $|A|=k-1$ and
$\cl(Ab)=\cl(Ab')$.  Choose $C$ of size $r-k$ such that
$AbC$ is a basis of $M$; then $Ab'C$ is also a basis.  In (1) for
$\nu\preceq\mu_M$ with lower set $A$ and upper set
$A\cup\{b,b'\}\cup C$, the only finite terms are those indexed by
$b$ and $b'$, since deleting an element of $C$ leaves the dependent
set $Abb'$.  Therefore $\nu(Ab)=\nu(Ab')$.  The case $k=0$ is
trivial.
\end{proof}

We may therefore write
\[
        \bar\nu(X)=\nu(B),\qquad \rk X=k,\quad \cl(B)=X.
\]

\section*{2. The opposition transform}

For a rank $k$ quotient $\nu\preceq\mu_M$ define a rank $r-k$
coordinate vector by
\[
 \nu^{\perp_M}(A)=
 \begin{cases}
   \bar\nu\bigl((\cl A)^\perp\bigr),
      & A\text{ independent in }M,\ |A|=r-k,\\
   \infty,&\text{otherwise.}
 \end{cases}                                                 \tag{4}
\]
This formula is independent of representatives and is visibly
involutive on flat coordinates.  The next lemma is the main local
calculation.

\begin{lemma}[elementary opposition]\label{lem:elementary}
Let $p\preceq q\preceq\mu_M$, with $\rk p=d$ and
$\rk q=d+1$.  Then $q^{\perp_M}$ and $p^{\perp_M}$ satisfy all
adjacent relations (2) of ranks $r-d-1$ and $r-d$.
\end{lemma}

\begin{proof}
If $d=r-1$, there are no relations to check.  Otherwise fix a set
$A$ of size $r-d-2$ and distinct $x,y,z\notin A$.  The three terms
to compare are
\[
 q^{\perp_M}(Ax)+p^{\perp_M}(Ayz),\quad
 q^{\perp_M}(Ay)+p^{\perp_M}(Axz),\quad
 q^{\perp_M}(Az)+p^{\perp_M}(Axy).                       \tag{5}
\]
If $A$ is dependent all three terms are $\infty$, so assume $A$ is
independent.  Put $U=\cl(A)$, $W=\cl(Axyz)$, and
$\delta=\rk W-\rk U$.

If $\delta\le1$, every two-element extension among $x,y,z$ is
dependent over $U$, so every $p^{\perp_M}$ term in (5) is $\infty$.

Assume $\delta=2$.  Let $G=W^\perp$ and $H=U^\perp$; then
$\rk G=d$ and $\rk H=d+2$.  Every finite term in (5) has the common
$p$-coordinate $\bar p(G)$.  For a finite term indexed by
$e\in\{x,y,z\}$, the other coordinate is
$\bar q(L_e)$, where
\[
        L_e=(\cl(Ae))^\perp
\]
is a rank $d+1$ flat in the rank-two interval $[G,H]$; terms for
which $e\in U$ or the complementary pair does not span $W/U$ are
$\infty$.  In a rank-two matroid on three labelled elements, either
the finite terms already contain two equal flats $L_e$, or the three
finite flats $L_x,L_y,L_z$ are distinct points of $[G,H]$.  In the
latter case choose a basis $S$ of $G$, representatives
$\ell_x,\ell_y,\ell_z$ in those three points, and extend
$S\ell_x\ell_y$ to a basis of $M$ by a set $C$ outside $H$.  The
incidence relation for $q\preceq\mu_M$ with lower set $S$ and upper
set $S\cup\{\ell_x,\ell_y,\ell_z\}\cup C$ has precisely the three
finite terms
\[
        \bar q(L_x),\qquad \bar q(L_y),\qquad \bar q(L_z),
\]
because deleting an element of $C$ leaves three points in a rank-two
interval.  Hence their minimum is attained at least twice.  Adding
the common summand $\bar p(G)$ proves (5).

It remains to treat $\delta=3$.  Then $G=W^\perp$ and $H=U^\perp$
have ranks $d-1$ and $d+2$.  In the rank-three modular interval
$[G,H]$, set
\[
 P_x=(\cl(Ayz))^\perp,\quad P_y=(\cl(Axz))^\perp,\quad
 P_z=(\cl(Axy))^\perp
\]
and
\[
 Q_x=(\cl(Ax))^\perp,\quad Q_y=(\cl(Ay))^\perp,\quad
 Q_z=(\cl(Az))^\perp .
\]
Since $x,y,z$ are independent over $U$, the flats
$\cl(Ayz),\cl(Axz),\cl(Axy)$ are the three planes spanned by pairs in
a rank-three frame over $U$.  By modularity their pairwise meets are
$\cl(Az),\cl(Ay),\cl(Ax)$, respectively.  Applying the
anti-isomorphism gives that $P_x,P_y,P_z$ are distinct atoms over
$G$ and
\[
        Q_x=P_y\vee P_z,\qquad
        Q_y=P_x\vee P_z,\qquad
        Q_z=P_x\vee P_y .
\]
Choose a basis $S$ of $G$ and elements
$a_x\in P_x\setminus G$, $a_y\in P_y\setminus G$,
$a_z\in P_z\setminus G$.  The three-term relation (2) for the
original elementary quotient $p\preceq q$, applied to
$S$ and $a_x,a_y,a_z$, is exactly
\[
 \min\{\bar p(P_x)+\bar q(Q_x),\,
        \bar p(P_y)+\bar q(Q_y),\,
        \bar p(P_z)+\bar q(Q_z)\}
\]
attained at least twice, which is (5).
\end{proof}

\section*{3. What the local calculation proves}

\begin{proposition}\label{prop:setinvolution}
The formula $\nu\mapsto \nu^{\perp_M}$ defines an involution of the
set underlying $\Dr(M)$.  It reverses every elementary inclusion:
if $\theta\preceq\nu\preceq\mu_M$ and
$\rk\nu=\rk\theta+1$, then
\[
        \nu^{\perp_M}\preceq\theta^{\perp_M}.
\]
Consequently it reverses every inclusion that can be refined by
elementary quotients.
\end{proposition}

\begin{proof}
Insert $\nu$ into a complete flag
$h_0\preceq h_1\preceq\cdots\preceq h_r=\mu_M$ using (3).
Lemma \ref{lem:elementary} applied to all adjacent pairs shows that
\[
        h_r^{\perp_M},h_{r-1}^{\perp_M},\ldots,h_0^{\perp_M}
\]
satisfies all adjacent three-term relations.  Its endpoints are the
trivial rank $0$ valuation and $\mu_M$: indeed $(\mu_M)^{\perp_M}$ is
rank $0$, while the transform of the rank $0$ valuation is the
trivial rank $r$ valuation supported on the bases of $M$.  By the
complete-flag criterion, the reversed sequence is a valuated complete
flag.  Thus $\nu^{\perp_M}\preceq\mu_M$, so the formula maps
$\Dr(M)$ to itself.  Since $(X^\perp)^\perp=X$ for flats, applying
(4) twice recovers $\nu$ up to a projective scalar.

If $\theta\preceq\nu$ is elementary, Lemma \ref{lem:elementary}
gives the adjacent incidence relations for
$\nu^{\perp_M},\theta^{\perp_M}$.  The preceding paragraph already
shows that both are valuated matroid quotients of $\mu_M$, so the
Brandt--Eur--Zhang incidence criterion gives
$\nu^{\perp_M}\preceq\theta^{\perp_M}$.  A chain of elementary
refinements gives the last assertion by iteration.
\end{proof}

\section*{4. Agreement with the flat involution}

Let $F$ be a flat of rank $f$ and let $q_F$ denote the rank $r-f$
valuated matroid of $M/F\oplus U_{0,F}$.  For an independent
$(r-f)$-set $B$ with $Z=\cl(B)$,
\[
        q_F(B)<\infty\quad\Longleftrightarrow\quad Z\vee F=\hat1.
                                                               \tag{6}
\]
For an independent $f$-set $A$ with $Y=\cl(A)$, (4) and (6) give
\[
 \Phi(q_F)(A)<\infty
 \quad\Longleftrightarrow\quad
        Y^\perp\vee F=\hat 1 .
\]
Because $\rk(Y^\perp)=r-f$ and $\rk F=f$, modularity makes this
equivalent to $Y^\perp\wedge F=\hat0$.  Applying the order-reversing
involution gives
\[
        Y\vee F^\perp=\hat1,
\]
which is precisely the condition that $A$ be a basis of
$M/F^\perp\oplus U_{0,F^\perp}$.  Dependent $A$ has value $\infty$ on
both sides.  Therefore the set involution of
Proposition \ref{prop:setinvolution} satisfies
\[
 \Phi\bigl(\Trop(M/F\oplus U_{0,F})\bigr)
 =
 \Trop(M/F^\perp\oplus U_{0,F^\perp}).
\]

\section*{Remaining open issues}

The construction above gives a canonical involution of the underlying
set of $\Dr(M)$, proves that it reverses all elementary quotients, and
therefore proves the desired order reversal for any inclusion
$\theta\preceq\nu$ that admits an elementary refinement
\[
        \theta=h_a\preceq h_{a+1}\preceq\cdots\preceq h_b=\nu .
\]
What is not yet proved here is that every inclusion in the relative
Dressian of such a modular matroid has such a refinement.  Equivalently,
one still needs either a relative partial-flag completion theorem in
this setting, or a direct proof that the full incidence relations
(1) for $\theta\preceq\nu$ are carried by (4) to the full incidence
relations for $\nu^{\perp_M}\preceq\theta^{\perp_M}$.  The local
rank-difference-one argument above was written to isolate exactly the
part of that direct proof that is currently under control.

\begin{thebibliography}{9}

\bibitem{BEZ}
M.~Brandt, C.~Eur and L.~Zhang,
\emph{Tropical flag varieties},
Advances in Mathematics \textbf{384} (2021), 107695.

\bibitem{DW}
A.~Dress and W.~Wenzel,
\emph{Valuated matroids},
Advances in Mathematics \textbf{93} (1992), 214--250.

\bibitem{JL}
M.~Jarra and O.~Lorscheid,
\emph{Flag matroids with coefficients},
Advances in Mathematics \textbf{436} (2024), 109396.

\bibitem{JLL0}
M.~Joswig, G.~Loho, D.~Luber and J.~A. Olarte,
\emph{Generalized permutahedra and positive flag Dressians},
International Mathematics Research Notices \textbf{2023} (2023),
16748--16777.

\bibitem{Murota}
K.~Murota,
\emph{Matroid valuation on independent sets},
Journal of Combinatorial Theory, Series B \textbf{69} (1997), 59--78.

\bibitem{MurotaBook}
K.~Murota,
\emph{Discrete Convex Analysis},
SIAM, Philadelphia, 2003.

\bibitem{Speyer}
D.~Speyer,
\emph{Tropical linear spaces},
SIAM Journal on Discrete Mathematics \textbf{22} (2008), 1527--1558.

\end{thebibliography}

\end{document}
